
	\newpage
	
% =========================================================	
	\section{Zahlendarstellung} \label{chap:zahlendarstellung}
% =========================================================
	
	\defi{
		\n 
		Zahlendarstellungen sind eine wichtige Grundlage der Numerik. Zunächst betrachten wir die Zahlenbasen, welche für uns relevant sind und wie wir sie übersetzen.
	%}
	%
	%
		\n
		Die wichtigsten Zahlenbasen sind für uns:
		\begin{itemize}
			\item Basis 2, binäre Zahlen genannt,
			\item Basis 10, die Dezimalzahlen und
			\item Basis 16 oder auch Hex-Zahlen.
		\end{itemize}
		Falls eine Darstellung nicht hinreichend gekennzeichnet ist, schreibe ich eine 2 in den Index bei binären Zahlen, bei Hex-Zahlen eine 0x vor die Zahl oder eine 16 in den Index. Falls die Darstellung durch einen Hex-Buchstaben trivialerweise einer Hex-Zahl entspricht, lasse ich eine Kennzeichnung weg.
%	}
%	
		\n 
		In der nächsten Tabelle lasse ich z.B. die Kennzeichnungen weg, da die erste Spalte die Basis der Zahlen angibt.
%	}
% =========================================================
%Zahlenbasen Beispiele (Tabelle)
	\begin{table}[H]
		\centering
			\begin{tabular}{lcccc}
				\toprule
				Basis   & \multicolumn{4}{c}{Beispiele} \\ \midrule
				Binär   & 01100  & 10001 & 11101 & 00011 \\
				Dezimal & 12     & 17    & 29    & 03    \\
				Hex     & 0C     & 11    & 1D    & 03    \\
				\bottomrule
			\end{tabular}%
		\caption{Beispiele für alle Basen.}
		\label{table:basen}
	\end{table}
% =========================================================
	\begin{table}[H]
		\centering
		\begin{tabular}{lccccccccccc}
			\toprule
			Ganzzahl & 1                   & 2                   & 4                   & 8                   & 16                  & 32                  & 64                  & 128                 & 256                 & 512                 & 1024                 \\ \midrule
			Potenzschreibweise & 2$^0$ & 2$^1$ & 2$^2$ & 2$^3$ & 2$^4$ & 2$^5$ & 2$^6$ & 2$^7$ & 2$^8$ & 2$^9$ & $2^{10}$ \\ 
			\bottomrule
		\end{tabular}
		\caption{Die wichtigsten Zweierpotenzen.}
		\label{table:2potenz}	
	\end{table}
% =========================================================
	\begin{table}[H]
		\centering
		\begin{tabular}{lcccccc}
			\toprule
			Hex-Buchstabe & A                   & B                   & C                   & D                   & E                  & F\\ \midrule
			Dezimalwert & 10 & 11 & 12 & 13 & 14 & 15\\ 
			\bottomrule
		\end{tabular}
		\caption{Wdh. Hex Buchstaben.}
		\label{table:hex}
	\end{table}
}

	\newpage
% =========================================================
	\subsection{Dezimal zu Binär}

		Du hast beispielweise die Zahl 12 und möchtest diese binär darstellen.\\
		
		\methode{
			\begin{itemize}
				\item Überlege dir, welche größte Zweierpotenz noch in deine Zahl reinpasst\\
				(Bei $12$ wäre das $8$, also $2^3$)
				\item Subtrahiere deine Zahl durch diese Potenz und schreibe eine \qu{$1$}.
				\item Geh jetzt jeweils durch alle kleineren Zweierpotenzen (bis $2^0$) und subtrahiere deine Zahl durch diese, wenn möglich \\
				(Solang die Subtraktion nicht kleiner null wird)
				\item Immer wenn man die Zahl nicht subtrahieren kann ($< 0$), schreibe eine \qu{$0$}, ansonsten eine \qu{$1$}.
			\end{itemize}
		}
		
		\sbreak	
		\bsp{(von oben)
			\begin{align*}
				12 - 2^{3} &= 4 \imp 1\\
				4 - 2^{2} &= 0 \imp 1\\
				0 - 2^{1} &< 0 \imp 0\\
				0 - 2^{0} &< 0 \imp 0
			\end{align*}
			Somit ergibt sich (von oben nach unten) \qu{$1100_2$}.
		}
				
		\sbreak
		\bsp{ Umwandlung von $37$:
			\begin{align*}
				37 - 2^{5} &= 5 \imp 1\\
				5 - 2^{4} &< 0 \imp 0\\
				5 - 2^{3} &< 0 \imp 0\\
				5 - 2^{2} &= 1 \imp 1\\
				1 - 2^{1} &< 0 \imp 0\\
				1 - 2^{0} &= 0 \imp 1
			\end{align*}
			Das Ergebnis ist \qu{$100101_2$}.
		}

\newpage
% =========================================================	
	\subsection{Binär zu Hex}
		\sbreak
	
		\methode{\n
			Wenn wir eine binäre Zahl in eine hexadezimale Zahl umwandeln wollen, ergeben vier binäre Ziffern jeweils eine Hex-Ziffer der Hex-Zahl.
			Also geht man durch die binäre Zahl von rechts nach links, trennt sie in Viererpakete auf und wandelt dann jedes Paket einzeln um.
		}
		
		\sbreak
		\bsp{ 
			\n
			$101010_2$ zu Hex umwandeln (Das wäre 42 als Dezimalzahl):
			\n
			Auftrennen in Viererpakete: \fat{0010|1010}
			\n 
			Die letzten (rechtesten) vier Ziffern ergeben im Dezimalsystem 10. Dies entspricht in Hex einem \qu{A} (siehe \refTable{table:hex}).
			Das andere Päkchen entspricht einer 2 und in Hex damit einer 2.
			\n 
			Die Hex Zahl ist somit \qu{2A}.
		}
		
		\sbreak
		\bsp{\n 
			Umwandlung von $110110011_2$: 
			\n
			Viererpakete: \fat{0001|1011|0011}
			\n
			Einzelne Umwandlungen: 
			\begin{align*}
				0001_2 &\imp 1& 1011_2 &\imp 11 (= \text{B})& 0011_2 \imp 3
			\end{align*}
			Die Hex Zahl ist somit "1B3"
		}
	
	\newpage
% =========================================================	
	\subsection{Dezimalbrüche zu Binär}
		
		\sbreak
		\methode{\n
			Führe schriftliches Dividieren der Brüche im binären System aus (das geht analog zum dezimalen).
			Das Ergebnis der Division ist auch das Ergebnis der Konvertierung.}
	
		\sbreak
		\bsp{\n
			$\frac{3}{4} = 011_2 : 100_2$
			\begin{table}[H]
				\label{table:dec>binary1}
				\begin{tabular}{llllll}
					0 	& 1 & 1 &   & \multicolumn{2}{l}{: $100_2$ = $0.11_2$} \\ \cline{1-4}
						& 1 & 1 & 0 & &  	\\ 
					- 	& 1 & 0 & 0 & &		\\ \cline{2-5}
						&   & 1 & 0 & 0 &   \\
						& - & 1 & 0 & 0 &   \\ \cline{3-5}
						& & & & 0 &
				\end{tabular}
			\end{table}
			\noindent Das Ergebnis ist also "$0.11_2$".
		}
		
		\sbreak
		\bsp{\n
			$\frac{2}{3} = 010_2 : 011_2$
			\begin{table}[H]
				\label{table:dec>binary2}
				\begin{tabular}{llllll}
					0 & 1 & 0 &   & \multicolumn{2}{l}{: $011_2$ = $0.101_2$} \\ \cline{1-4}
					& 1 & 0 & 0 & \\
					- & 0 & 1 & 1 & \\ \cline{2-4}
					& 0 & 0 & 1 & 0 & (Periodisch)
				\end{tabular}
			\end{table}
			\noindent Hier haben wir den Fall, dass sich die Rechnung ständig wiederholt,
			als nächstes würden wieder eine Null und eine Eins kommen usw.
			\n 
			Das Ergebnis ist also: 
			\begin{align*}
				0.1\overline{01}_2
			\end{align*}
		}
